% Generated draft: requires mathematical and source review.
\documentclass[11pt]{article}
\usepackage[margin=25mm]{geometry}
\usepackage{fontspec}
\setmainfont{Latin Modern Roman}
\usepackage{amsmath,amssymb,mathtools,unicode-math}
\setmathfont{Latin Modern Math}
\usepackage{enumitem,needspace,xurl,hyperref}
\setlength{\parindent}{0pt}
\setlength{\parskip}{5pt}
\setlength{\emergencystretch}{3em}
\newcommand{\N}{\mathbb{N}}
\newcommand{\Z}{\mathbb{Z}}
\newcommand{\Q}{\mathbb{Q}}
\newcommand{\R}{\mathbb{R}}
\newcommand{\C}{\mathbb{C}}
\newcommand{\F}{\mathbb{F}}
\newcommand{\A}{\mathbb{A}}
\newcommand{\PP}{\mathbb{P}}
\newcommand{\E}{\mathbb{E}}
\newcommand{\eps}{\varepsilon}
\newcommand{\dd}{\,\mathrm{d}}
\begin{document}
\section*{UnsolvedMath ID: 800768}
\textbf{Source problem code:} catalogue entry
\subsection{Definitions and mathematical statement}
Let \(X\) be a compact K\"ahler manifold of complex dimension \(n\), equipped with a K\"ahler metric \(g\). Assume that \(g\) has semi-positive holomorphic sectional curvature, meaning that its holomorphic sectional curvature is nonnegative in every nonzero holomorphic tangent direction. Impose the additional source hypothesis
\[
n_{\mathrm{tf}}(X,g)=n.
\]
The supplied record does not define \(n_{\mathrm{tf}}(X,g)\). This hypothesis is therefore retained as written; its intended definition must be supplied before the class of manifolds under consideration is fully specified.

For each positive integer \(q\), let \(\Omega_X^q\) denote the sheaf of holomorphic differential \(q\)-forms, and set
\[
h^0(X,\Omega_X^q)=\dim_{\C}H^0(X,\Omega_X^q).
\]
The source asks the following two questions for every pair \((X,g)\) satisfying the stated hypotheses:
\begin{enumerate}
\item Must
\[
H^0(X,\Omega_X^q)=0
\qquad\text{for every integer }q>0?
\]
\item Must \(X\) be projective, that is, admit a holomorphic embedding into \(\PP^N_{\C}\) for some finite integer \(N\geq 0\)?
\end{enumerate}
For \(q>n\), the sheaf \(\Omega_X^q\) is zero, so the substantive vanishing question concerns \(1\leq q\leq n\). The record labels the problem open; that label is supplied source metadata, not an independently verified assessment.
\subsection{Short English statement}
Matsumura asks whether a compact K\"ahler manifold with semi-positive holomorphic sectional curvature and \(n_{\mathrm{tf}}(X,g)=\dim_{\C}X\) must have no nonzero holomorphic forms of positive degree, and whether it must be projective. The invariant \(n_{\mathrm{tf}}\) is undefined in the supplied record, so the hypotheses require clarification.
\subsection{Sources}
\begin{itemize}
\item Local UnsolvedMath registry, supplied record with permanent ID \(800768\), exact source code \texttt{catalogue entry}, in the AMR Open Problem Lists collection. The record attributes the problem to Matsumura, with proposed year 2022, and labels it open.
\item Matsumura, \emph{Open Problems on Structure of Positively Curved Projective Varieties} (2022), Problem 3. Supplied source URL: \url{https://arxiv.org/abs/2211.09153}. The record reports access on 2026-10-16 and extraction method \texttt{arxiv\_latex}; these are supplied provenance statements, not claims of independent access or verification.
\item The supplied statement additionally cites \texttt{[Mat18b, Problem 5.1]}. No bibliographic entry identifying \texttt{Mat18b} is supplied, so this attribution cannot be resolved from the record alone.
\end{itemize}
\end{document}
